\documentclass[a4paper,12pt]{article}
\usepackage[utf8]{inputenc}
\usepackage[T1]{fontenc}
\usepackage[ngerman]{babel}
\usepackage{amsmath}
\usepackage{amssymb}
\usepackage{enumitem}
\usepackage{geometry}
\geometry{a4paper, margin=2.5cm}

\title{Einsendeaufgaben -- Lektion 3\\[0.5em]
\large Modul 61111: Mathematische Grundlagen}
\author{}
\date{}

\begin{document}
\maketitle

\section*{Aufgabe 3.3}

\begin{enumerate}[label=\alph*)]

    \item
    \[
        f_1(v+w)
        = \begin{pmatrix}v_1+w_1-v_2-w_2\\2v_1+2w_1\\3v_2+3w_2\end{pmatrix}
        = \begin{pmatrix}v_1-v_2\\2v_1\\3v_2\end{pmatrix}
        + \begin{pmatrix}w_1-w_2\\2w_1\\3w_2\end{pmatrix}
        = f_1(v) + f_1(w)
    \]

    \[
        f_1(\lambda v)
        = \begin{pmatrix}\lambda x - \lambda y\\2\lambda x\\3\lambda y\end{pmatrix}
        = \lambda \begin{pmatrix}x-y\\2x\\3y\end{pmatrix}
        = \lambda f_1(v)
    \]

    $\Rightarrow f_1$ ist eine lineare Abbildung.

    \item keine lineare Abbildung, da
    \[
        f_2\!\left(\lambda\begin{pmatrix}x\\y\end{pmatrix}\right)
        = \begin{pmatrix}\lambda x + \lambda^2 y^2\\2\lambda x\\3\lambda y - 1\end{pmatrix}
        \neq
        \begin{pmatrix}\lambda x + \lambda y^2\\2\lambda x\\3\lambda y - \lambda\end{pmatrix}
        = \lambda f_2\!\left(\begin{pmatrix}x\\y\end{pmatrix}\right)
    \]

    \item keine lineare Abbildung, da
    \[
        f_3\!\left(\lambda\begin{pmatrix}x\\y\end{pmatrix}\right)
        = \begin{pmatrix}\lambda(x-y)\\2\lambda^2 xy\\3\lambda y\end{pmatrix}
        \neq
        \begin{pmatrix}\lambda(x-y)\\2\lambda xy\\3\lambda y\end{pmatrix}
        = \lambda f_3\!\left(\begin{pmatrix}x\\y\end{pmatrix}\right)
    \]

\end{enumerate}

\end{document}
